% ======================================================================
% Chapter 12: The Renormalization Group
% (book pages 393-438; PDF pages 416-461)
% ======================================================================
\chapter{The Renormalization Group}

In the past two chapters, our main goal has been to determine when, and how,
the cancellation of ultraviolet divergences in quantum field theory takes
place.  We have seen that, in a large class of field theories, the divergences
appear only in the values of a few parameters: the bare masses and coupling
constants, or, in renormalized perturbation theory, the counterterms.  Aside
from the shift in these parameters, virtual particles with very large momenta
have no effect on computations in these theories.

The cancellation of ultraviolet divergences is essential if a theory is to yield
quantitative physical predictions.  But, at a deep level, the fact that
high-momentum virtual quanta can have so little effect on a theory is quite
surprising.  One of the essential features of quantum field theory is
locality, that is, the fact that fields at different spacetime points are
independent degrees of freedom with independent quantum fluctuations.  The
quantum fluctuations at arbitrarily short distances appear in Feynman diagram
computations as virtual quanta with arbitrarily high momenta.  In a
renormalizable theory, the loop integrals over virtual-particle momenta are
always dominated by values comparable to the finite external particle
momenta.  But why?  It is not easy to understand how the quantum fluctuations
associated with extremely short distances can be so innocuous as to affect a
theory only through the values of a few of its parameters.

This chapter begins with a physical picture, due to Kenneth Wilson, that
explains this unusual and counterintuitive simplification.  This picture
generalizes the idea of the distance- or scale-dependent electric charge,
introduced at the end of Chapter~7, and suggests that all of the parameters of
a renormalizable field theory can usefully be thought of as scale-dependent
entities.  We will see that this scale dependence is described by simple
differential equations, called \emph{renormalization group equations}.  The
solutions of these equations will lead to physical predictions of a completely
new type: predictions that, under certain circumstances, the correlation
functions of a quantum field exhibit unusual but computable scaling laws as a
function of their coordinates.

\section{Wilson's Approach to Renormalization Theory}

Wilson's method is based on the functional integral approach to field theory,
in which the degrees of freedom of a quantum field are variables of
integration.  In this approach, one can study the origin of ultraviolet
divergences by isolating the dependence of the functional integral on the
short-distance degrees of freedom of the field.%
\footnote{Wilson's ideas are reviewed in K.~G.~Wilson and J.~Kogut,
\emph{Phys.\ Repts.} \textbf{12C}, 75 (1974).}  In this section, we will
illustrate this idea in the simplest example of $\phi^4$ theory.

To make our analysis more concrete, we will drop the elegant but somewhat
mysterious method of dimensional regularization in this section and instead
use a sharp momentum cutoff.  Since we will be working here only in $\phi^4$
theory, we will not be concerned that this cutoff makes it difficult to
satisfy Ward identities.  Wilson's analysis can be adapted to QED and other
situations where this subtlety is important, but the case of $\phi^4$ theory
is sufficient to give us the basic qualitative results of this approach.

In Section~9.2, we constructed the Green's functions of $\phi^4$ theory in
terms of a functional integral representation of the generating functional
$Z[J]$.  The basic integration variables are the Fourier components of the
field $\phi(k)$, so $Z[J]$ is given concretely by the expression
\begin{equation}
  Z[J] = \int \mathcal{D}\phi\ e^{i[S+J\phi]}
  = \Bigl( \prod_{k_n} \int d\phi(k_n) \Bigr)\ e^{i[S+J\phi]} .
  \label{eq:12.1}
\end{equation}
To impose a sharp ultraviolet cutoff $\Lambda$, we restrict the number of the
integration variables displayed in \eqref{eq:12.1}.  That is, we integrate
only over $\phi(k)$ with $|k| < \Lambda$, and set $\phi(k) = 0$ for
$|k| > \Lambda$.

This modification of the functional integral suggests a method for assessing
the influence of the quantum fluctuations at very short distances or very
large momenta.  In the functional integral representation, these fluctuations
are represented by the integrals over the Fourier components of $\phi$ with
momenta near the cutoff.  Why not explicitly perform the integrals over these
variables?  Then we can compare the result to the original functional
integral, and determine precisely the influence of these high-momentum modes
on the physical predictions of the theory.

Before beginning this analysis, though, we must introduce one modification.
At first sight, it seems most natural to define the ultraviolet cutoff in
Minkowski space.  However, a cutoff $k^2 < \Lambda^2$ is not completely
effective in controlling large momenta, since in lightlike directions the
components of $k$ can be very large while $k^2$ remains small.  We will
therefore consider the cutoff to be imposed on the Euclidean momenta obtained
after Wick rotation.  Equivalently, we consider the Euclidean form of the
functional integral, presented in Section~9.3, and restrict its variables
$\phi(k)$, with $k$ Euclidean, to $|k| < \Lambda$.

The transition to Euclidean space also brings us closer to the connection
between renormalization theory and statistical mechanics advertised in
Chapter~8.  As we saw in Section~9.3, the Euclidean functional integral for
$\phi^4$ theory has precisely the same form as the continuum description of
the statistical mechanics of a magnet.  The field $\phi(x)$ is interpreted as
the fluctuating spin field $s(x)$.  A real magnet is built of atoms, and the
atomic spacing provides a physical cutoff, a shortest distance over which
fluctuations can take place.  The cut-off functional integral models the
effects of this atomic size in a crude way.

By pursuing this analogy, we can derive some physical intuition about the
effects of the ultraviolet cutoff in a field theory.  In a magnet, it is quite
easy to visualize statistical fluctuations of the spins at the atomic scale.
In fact, for values of the temperature away from any critical points, the
statistical fluctuations are restricted to this scale; over distances of tens
of atomic spacings, the magnet already shows its homogeneous macroscopic
behavior.  We have seen in Chapter~8 that we can approximate the correlation
function of the spin field by the propagator of a Euclidean $\phi^4$ theory.
In this approximation,
\begin{equation}
  \langle s(x)s(0)\rangle \sim \int \frac{d^dk}{(2\pi)^d}\
    \frac{e^{ik\cdot x}}{k^2+m^2} .
  \label{eq:12.2}
\end{equation}
As long as the temperature is far from the critical temperature, the size of
the ``mass'' $m$ is determined by the one natural scale in the problem, the
atomic spacing.  Thus, we expect $m \sim \Lambda$.  In our field theory
calculations, we were specifically interested in the situation where
$m \ll \Lambda$, and we adjusted the parameters of the theory to satisfy this
condition.  In describing a magnet, it appears that no such adjustment is
called for.

However, we saw in Chapter~8 that there is one circumstance in which the
correlations of the spin field are much longer than the atomic spacing, so
that, indeed, $m \ll \Lambda$.  When the spin system begins to magnetize, just
in the vicinity of the critical point, the spins become correlated over
arbitrarily long distances as the fluctuating spins attempt to choose their
eventual direction of magnetization.  To study these long-range correlations
in a magnet, one must carefully adjust the temperature to bring the system
into the vicinity of the phase transition.  In the same way, we can imagine
making a fine adjustment of the parameter $m$ of $\phi^4$ theory to bring the
quantum field theory into a region of parameters where we do find
correlations of the field $\phi(x)$ over distances much larger than
$1/\Lambda$.

\subsection{Integrating Over a Single Momentum Shell}

With this introduction, we will now carry out the integration over the
high-momentum degrees of freedom of $\phi$.  We begin by writing the
functional integral \eqref{eq:12.1} more explicitly for the case of $\phi^4$
theory.  We apply the cutoff prescription described earlier, and set $J = 0$
for simplicity.  Then
\begin{equation}
  Z = \int [\mathcal{D}\phi]_\Lambda\ \exp\Bigl( -\int d^dx\ \Bigl[
    \tfrac{1}{2}(\partial_\mu\phi)^2 + \tfrac{1}{2}m^2\phi^2
    + \frac{\lambda}{4!}\phi^4 \Bigr] \Bigr),
  \label{eq:12.3}
\end{equation}
where
\begin{equation}
  [\mathcal{D}\phi]_\Lambda = \prod_{|k|<\Lambda} d\phi(k) .
  \label{eq:12.4}
\end{equation}
In the Lagrangian of Eq.~\eqref{eq:12.3}, $m$ and $\lambda$ are the bare
parameters, and so there are no counterterms.  As in our study of the
superficial degree of divergence, it will be useful to carry out this analysis
in an arbitrary spacetime dimension $d$.

We now divide the integration variables $\phi(k)$ into two groups.  Choose a
fraction $b < 1$.  The variables $\phi(k)$ with $b\Lambda < |k| < \Lambda$ are
the high-momentum degrees of freedom that we will integrate over.  To label
these degrees of freedom, let us define
\begin{equation}
  \hat{\phi}(k) = \begin{cases}
    \phi(k) & \text{for } b\Lambda < |k| < \Lambda;\\
    0 & \text{otherwise}.
  \end{cases}
\end{equation}
Next, let us define a new $\phi(k)$, which is identical to the old for
$|k| < b\Lambda$ and zero for $|k| > b\Lambda$.  Then we can replace the old
$\phi$ in the Lagrangian with $\phi + \hat{\phi}$, and rewrite Eq.
\eqref{eq:12.3} as
\begin{equation}
  \begin{aligned}
    Z &= \int \mathcal{D}\phi\ \int \mathcal{D}\hat{\phi}\ \exp\Bigl(
      -\int d^dx\ \Bigl[ \tfrac{1}{2}(\partial_\mu(\phi+\hat{\phi}))^2
        + \tfrac{1}{2}m^2(\phi+\hat{\phi})^2 + \frac{\lambda}{4!}
          (\phi+\hat{\phi})^4 \Bigr] \Bigr)\\
    &= \int \mathcal{D}\phi\ e^{-\int d^dx\ \mathcal{L}(\phi)}\ \int
      \mathcal{D}\hat{\phi}\ \exp\Bigl( -\int d^dx\ \Bigl[
        \tfrac{1}{2}(\partial_\mu\hat{\phi})^2 + \tfrac{1}{2}m^2\hat{\phi}^2
        + \frac{\lambda}{4!}\hat{\phi}^4 + \frac{\lambda}{3!}\hat{\phi}^3\phi
        + \frac{\lambda}{4}\hat{\phi}^2\phi^2 \Bigr] \Bigr) .
  \end{aligned}
  \label{eq:12.5}
\end{equation}
In the final expression we have gathered all terms independent of
$\hat{\phi}$ into $\mathcal{L}(\phi)$.  Note that quadratic terms of the form
$\phi\hat{\phi}$ vanish, since Fourier components of different wavelengths are
orthogonal.

The next few paragraphs will explain how to perform the integral over
$\hat{\phi}$.  This integration will transform \eqref{eq:12.5} into an
expression of the form
\begin{equation}
  Z = \int [\mathcal{D}\phi]_{b\Lambda}\ \exp\Bigl( -\int d^dx\
    \mathcal{L}_{eff}(\phi) \Bigr) ,
  \label{eq:12.6}
\end{equation}
where $\mathcal{L}_{eff}(\phi)$ involves only the Fourier components
$\phi(k)$ with $|k| < b\Lambda$.  We will see that $\mathcal{L}_{eff}(\phi) =
\mathcal{L}(\phi)$ plus corrections proportional to powers of $\Lambda$.
These correction terms compensate for the removal of the large-$k$ Fourier
components $\hat{\phi}$, by supplying the interactions among the remaining
$\phi(k)$ that were previously mediated by fluctuations of the
$\hat{\phi}$.

To carry out the integrals over the $\hat{\phi}(k)$, we use the same method
that we applied in Section~9.2 to derive Feynman rules.  In fact, we will see
below that the new terms in $\mathcal{L}_{eff}$ can be written in a
diagrammatic form.  In this analysis, we treat the quartic terms in
\eqref{eq:12.5}, all proportional to $\lambda$, as perturbations.  Since we
are mainly interested in the situation $m^2 \ll \Lambda^2$, we will also treat
the mass term $\tfrac{1}{2}m^2\hat{\phi}^2$ as a perturbation.  Then the
leading-order term in the portion of the Lagrangian involving $\hat{\phi}$ is
\begin{equation}
  \int \mathcal{L}_0 = \int_{b\Lambda<|k|<\Lambda} \frac{d^dk}{(2\pi)^d}\
    \frac{1}{2}k^2\,|\hat{\phi}(k)|^2 .
  \label{eq:12.7}
\end{equation}
This term leads to a propagator
\begin{equation}
  \langle\hat{\phi}(k)\hat{\phi}(p)\rangle = (2\pi)^d\,\delta(k+p)\,
    \frac{1}{k^2}\,g(k) ,
  \label{eq:12.8}
\end{equation}
where
\begin{equation}
  g(k) = \begin{cases}
    1 & \text{if } b\Lambda < |k| < \Lambda;\\
    0 & \text{otherwise}.
  \end{cases}
  \label{eq:12.9}
\end{equation}
We will regard the remaining $\hat{\phi}$ terms in Eq.~\eqref{eq:12.5} as
perturbations, and expand the exponential.  The various contributions from
these perturbations can be evaluated by using Wick's theorem with
\eqref{eq:12.8} as the propagator.

First consider the term that results from expanding to one power of the
$\hat{\phi}^2\phi^2$ term in the exponent of \eqref{eq:12.5}.  We find
\begin{equation}
  -\frac{\lambda}{4}\int d^dx\ \phi^2\hat{\phi}^2 = -\frac{1}{2}\int
    \frac{d^dk}{(2\pi)^d}\ \Delta m^2\,\phi(k)\phi(-k) ,
  \label{eq:12.10}
\end{equation}
where the coefficient $\Delta m^2$ is the result of contracting the two
$\hat{\phi}$ fields:
\begin{equation}
  \begin{aligned}
    \Delta m^2 &= \frac{\lambda}{2}\ \int_{b\Lambda<|k|<\Lambda}
      \frac{d^dk}{(2\pi)^d}\ \frac{1}{k^2}\\
    &= \frac{\lambda}{2}\ \frac{1}{(4\pi)^{d/2}}\ \frac{\Lambda^{d-2}}{\Gamma(d/2)}
      \ \frac{1-b^{d-2}}{d-2} .
  \end{aligned}
  \label{eq:12.11}
\end{equation}
The term \eqref{eq:12.10} could just as well have arisen from an expansion of
the exponential
\begin{equation}
  \exp\Bigl( -\int d^dx\ \frac{1}{2}\Delta m^2\,\phi^2 + \cdots \Bigr) .
  \label{eq:12.12}
\end{equation}
We will soon see that the rest of the perturbation series also organizes
itself into this form.  The coefficient $\Delta m^2$ therefore gives a positive
correction to the $m^2$ term in $\mathcal{L}$.

The higher orders of the perturbation theory in the correction terms can be
worked out in a similar way.  As in our derivation of the standard
perturbation theory for $\phi^4$ theory, it is useful to adopt a diagrammatic
notation.  Represent the propagator \eqref{eq:12.8} by a double line.  This
propagator will connect pairs of fields $\hat{\phi}$ from the various quartic
interactions.  Represent the fields $\phi$ in these interactions, which are
not integrated over, as single external lines.

The leading correction to the $\phi^4$ coupling constant comes from the
diagram with two internal $\hat{\phi}$ lines and four external $\phi$ lines:
\begin{equation}
  \Delta\lambda = -\frac{3\lambda^2}{2}\ \int_{b\Lambda<|k|<\Lambda}
    \frac{d^dk}{(2\pi)^d}\ \frac{1}{k^2}\ \frac{1}{k^2}
  = -\frac{3\lambda^2}{2}\ \frac{\Lambda^{d-4}}{(4\pi)^{d/2}}\ \frac{1-b^{d-4}}
    {d-4}\ \cdots ,
  \label{eq:12.15}
\end{equation}
up to finite terms.  Corrections to the field-strength term $(\partial_\mu\phi)^2$
and the coefficient of $(\partial_\mu\phi)^4$ come from diagrams with derivative
couplings and are suppressed by additional powers of $\lambda$.

We can now summarize the result of integrating over the high-momentum shell.
The effective Lagrangian has the form
\begin{equation}
  \int d^dx\ \mathcal{L}_{eff} = \int d^dx\ \Bigl[
    \tfrac{1}{2}(1+\Delta Z)(\partial_\mu\phi)^2 + \tfrac{1}{2}
      (m^2+\Delta m^2)\phi^2 + \frac{1}{4!}(\lambda+\Delta\lambda)\phi^4
      + \cdots \Bigr] .
  \label{eq:12.18}
\end{equation}
To compare this to the original Lagrangian, we must rescale the coordinates
and the field so that the kinetic term returns to its canonical form.  Define
rescaled coordinates
\begin{equation}
  x' = bx,\qquad \partial_\mu' = \frac{1}{b}\partial_\mu ,
  \label{eq:12.19}
\end{equation}
and rescale the field as
\begin{equation}
  \phi' = \bigl[ b^{2-d}(1+\Delta Z) \bigr]^{1/2}\phi .
  \label{eq:12.22}
\end{equation}
After this rescaling, the unperturbed action returns to its initial form,
while the various perturbations undergo a transformation.  Writing the result
in the form of a new Lagrangian, the new parameters of the Lagrangian are
\begin{equation}
  \begin{aligned}
    m'^2 &= (m^2+\Delta m^2)(1+\Delta Z)^{-1}\,b^{-2};\\
    \lambda' &= (\lambda+\Delta\lambda)(1+\Delta Z)^{-2}\,b^{d-4};\\
    C' &= (C+\Delta C)(1+\Delta Z)^{-2}\,b^d;\\
    D' &= (D+\Delta D)(1+\Delta Z)^{-3}\,b^{2d-6} ,
  \end{aligned}
  \label{eq:12.24}
\end{equation}
and so on.  (The original Lagrangian had $C = D = 0$, but the same equations
would apply if the initial values of $C$ and $D$ were nonzero.)  All of the
corrections, $\Delta m^2$, $\Delta\lambda$, and so on, arise from diagrams and
thus are small compared to the leading terms if perturbation theory is
justified.

By combining the operation of integrating out high-momentum degrees of
freedom with the rescaling \eqref{eq:12.19}, we have rewritten this operation
as a transformation of the Lagrangian.  Continuing this procedure, we could
integrate over another shell of momentum space and transform the Lagrangian
further.  Successive integrations produce further iterations of the
transformation \eqref{eq:12.24}.  If we take the parameter $b$ to be close to
1, so that the shells of momentum space are infinitesimally thin, the
transformation becomes a continuous one.  We can then describe the result of
integrating over the high-momentum degrees of freedom of a field theory as a
trajectory or a flow in the space of all possible Lagrangians.

For historical reasons, these continuously generated transformations of
Lagrangians are referred to as the \emph{renormalization group}.  They do not
form a group in the formal sense, because the operation of integrating out
degrees of freedom is not invertible.  On the other hand, they are most
certainly connected to renormalization, as we will now see.

Imagine that we wish to compute a correlation function of fields whose momenta
$p_i$ are all much less than $\Lambda$.  We could compute this correlation
function perturbatively using either the original Lagrangian $\mathcal{L}$, or
the effective Lagrangian $\mathcal{L}_{eff}$ obtained after integrating over
all momentum shells down to the scale of the external momenta $p_i$.  Both
procedures must ultimately yield the same result.  But in the first case, the
effects of high-momentum fluctuations of the field do not show up until we
compute loop diagrams.  In the second case, these effects have already been
absorbed into the new coupling constants ($m'$, $\lambda'$, etc.), so their
influence can be seen directly from the Lagrangian.  In the first procedure,
the large shifts from the original (bare) parameters to the values
appropriate to low-momentum processes appear suddenly in one-loop diagrams,
and seem to invalidate the use of perturbation theory.  In the second
approach, these corrections are introduced slowly and systematically.  A
perturbative treatment is valid at every step as long as the effective
coupling constants such as $\lambda'$ remain small.

However, the parameters of the effective Lagrangian may be very different from
those of the original Lagrangian, since we must iterate the transformation
\eqref{eq:12.24} many times to get from the large momentum $\Lambda$ down to
the momentum scale of typical experiments.  Let us therefore look more
closely at how the Lagrangian tends to vary under the renormalization group
transformations.

The simplest case to consider is a Lagrangian in the vicinity of the point
$m^2 = \lambda = C = D = \cdots = 0$, where all the perturbations vanish.  We
have defined our transformation so that this point is left unchanged; we say
that the free-field Lagrangian
\begin{equation}
  \mathcal{L}_0 = \tfrac{1}{2}(\partial_\mu\phi)^2
  \label{eq:12.25}
\end{equation}
is a \emph{fixed point} of the renormalization group transformation.

In the vicinity of $\mathcal{L}_0$, we can ignore the terms $\Delta m^2$,
$\Delta\lambda$, etc., in the iteration equations \eqref{eq:12.24} and keep
only those terms that are linear in the perturbations.  This gives an
especially simple transformation law:
\begin{equation}
  m'^2 = m^2b^{-2},\qquad \lambda' = \lambda b^{d-4},\qquad C' = Cb^d,\qquad
  D' = Db^{2d-6},\qquad \text{etc.}
  \label{eq:12.26}
\end{equation}
Since $b < 1$, those parameters that are multiplied by negative powers of $b$
grow, while those that are multiplied by positive powers of $b$ decay.  If the
Lagrangian contains growing coefficients, these will eventually carry it away
from $\mathcal{L}_0$.

It is conventional to speak of the various terms in the effective Lagrangian
as a set of local operators that can be added as perturbations to
$\mathcal{L}_0$.  We call the operators whose coefficients grow during the
recursion procedure \emph{relevant} operators.  The coefficients that die
away are associated with \emph{irrelevant} operators.  For example, the
scalar field mass operator $\phi^2$ is always relevant, while the $\phi^4$
operator is relevant if $d < 4$.  If the coefficient of some operator is
multiplied by $b^0$ (for example, the operator $\phi^4$ in $d = 4$), we call
this operator \emph{marginal}; to find out whether its coefficient grows or
decays, we must include the effect of higher-order corrections.

In general, an operator with $N$ powers of $\phi$ and $M$ derivatives has a
coefficient that transforms as
\begin{equation}
  C'_{N,M} = b^{N(d/2-1)+M-d}\,C_{N,M} .
  \label{eq:12.27}
\end{equation}
Notice that the coefficient is just $(d_{N,M}-d)$, where $d_{N,M}$ is the
mass dimension of the operator as computed at the end of Section~10.1.  In
other words, relevant and marginal operators about the free theory
$\mathcal{L}_0$ correspond precisely to super-renormalizable and
renormalizable interaction terms in the power-counting analysis of
Section~10.1.

We can also understand the evolution of coefficients near the free-field fixed
point using straightforward dimensional analysis.  An operator with mass
dimension $d_i$ has a coefficient with dimension $(\mathrm{mass})^{d-d_i}$.
The natural order of magnitude for this mass is the cutoff $\Lambda$.  Thus,
if $d_i < d$, the perturbation is increasingly important at low momenta.  On
the other hand, if $d_i > d$, the relative size of this term decreases as
$(p/\Lambda)^{d_i-d}$ as the momentum $p\to 0$; thus the term is truly
irrelevant.

We have now shown that, at least in the vicinity of the zero-coupling fixed
point, an arbitrarily complicated Lagrangian at the scale of the cutoff
degenerates to a Lagrangian containing only a finite number of renormalizable
interactions.  It is instructive to compare this result with the conclusions
of Chapter~10.  There we took the philosophy that the cutoff $\Lambda$ should
be disposed of by taking the limit $\Lambda\to\infty$ as quickly as possible.
We found that this limit gives well-defined predictions only if the
Lagrangian contains no parameters with negative mass dimension.  From this
viewpoint, it seemed exceedingly fortunate that QED, for example, contained
no such parameters, since otherwise this theory would not yield well-defined
predictions.

Wilson's analysis takes just the opposite point of view, that any quantum
field theory is defined fundamentally with a cutoff $\Lambda$ that has some
physical significance.  In statistical mechanical applications, this momentum
scale is the inverse atomic spacing.  In QED and other quantum field theories
appropriate to elementary particle physics, the cutoff would have to be
associated with some fundamental graininess of spacetime, perhaps a result of
quantum fluctuations in gravity.  We discuss some speculations on the nature
of this cutoff in the Epilogue.  But whatever this scale is, it lies far
beyond the reach of present-day experiments.  The argument we have just given
shows that this circumstance explains the renormalizability of QED and other
quantum field theories of particle interactions.  Whatever the Lagrangian of
QED was at its fundamental scale, as long as its couplings are sufficiently
weak, it must be described at the energies of our experiments by a
renormalizable effective Lagrangian.

\begin{figure}[t]
  \centering
  \includegraphics[width=0.6\textwidth]{fig12.1.png}
  \caption{Renormalization group flows near the free-field fixed point in
    scalar field theory: (a) $d > 4$; (b) $d = 4$.}
  \label{fig:12.1}
\end{figure}

On the other hand, we should emphasize that these simple conclusions can be
altered by sufficiently strong field theory interactions.  Away from the
free-field fixed point, the simple transformation laws \eqref{eq:12.26}
receive corrections proportional to higher powers of the coupling constants.
If these corrections are large enough, they can halt or reverse the
renormalization group flow.  They could even create new fixed points, which
would give new types of $\Lambda\to\infty$ limits.

To illustrate the possible influences of interactions in a relatively simple
context, let us discuss the renormalization group flows near $\mathcal{L}_0$
for the specific case of $\phi^4$ theory.  It is instructive to consider the
three cases $d > 4$, $d = 4$, and $d < 4$ in turn.  When $d > 4$, the only
relevant operator is the scalar field mass term.  Then the renormalization
group flows near $\mathcal{L}_0$ have the form shown in
Fig.~\ref{fig:12.1}(a).  The $\phi^4$ interaction and possible higher-order
interactions die away, while the mass term increases in importance.

In previous chapters, we have always discussed $\phi^4$ theory in the limit in
which the mass is small compared to the cutoff.  Let us take a moment to
rewrite this condition in the language of renormalization group flows.  In
the course of the flow, the effective mass term $m'^2$ becomes large and
eventually comes to equal the current cutoff.  For example, near the
free-field fixed point, after $n$ iterations, $m'^2 = m^2b^{-2n}$, and
eventually there is an $n$ such that $m'^2 \sim \Lambda^2$.  At this point, we
have integrated out the entire momentum region between the original
$\Lambda$ and the effective mass of the scalar field.  The mass term then
suppresses the remaining quantum fluctuations.  In general, the criterion
that the scalar field mass is small compared to the cutoff is equivalent to
the statement that $m'^2 \sim \Lambda^2$ only after a large number of
iterations of the renormalization group transformation.

This criterion is met whenever the initial conditions for the renormalization
group flow are adjusted so that the trajectory passes very close to a fixed
point.  In principle, the flow could begin far away, along the direction of an
irrelevant operator.  The original value of $m^2$ need not be particularly
small, as long as this original value is canceled by corrections arising from
the diagrammatic contributions to $\mathcal{L}_{eff}$.  Thus we could imagine
constructing a scalar field theory in $d > 4$ by writing a complicated
nonlinear Lagrangian, but adjusting the original $m^2$ so the trajectory that
begins at this Lagrangian eventually passes close to the free-field fixed
point $\mathcal{L}_0$.  In this case, the effective theory at momenta small
compared to the cutoff should be extremely simple: It will be a free field
theory with negligible nonlinear interaction.  As will be discussed in the
next chapter, this remarkable prediction has been verified in mathematical
models of magnetic systems in more than four dimensions: Even though the
original model is highly nonlinear, the correlation function of spins near
the phase transition has the free-field form given by the higher-dimensional
analogue of Eq.~\eqref{eq:12.2}.

Next consider the case $d = 4$.  For this case, Eq.~\eqref{eq:12.26} does not
give enough information to tell us whether the $\phi^4$ interaction is
important or unimportant at large distances.  So we must go back to the
complete transformation law \eqref{eq:12.24}.  The leading contribution to
$\Delta\lambda$ is given by Eq.~\eqref{eq:12.15}.  The leading contribution
to $\Delta Z$ is of order $\lambda^2$ and can be neglected.  (This is just
what happened with the first correction to $\delta_Z$ in Section~10.2.)  Thus
we find the transformation
\begin{equation}
  \lambda' = \lambda - \frac{3\lambda^2}{16\pi^2}\log(1/b) .
  \label{eq:12.28}
\end{equation}
This says that $\lambda$ slowly decreases as we integrate out high-momentum
degrees of freedom.

\begin{figure}[t]
  \centering
  \includegraphics[width=0.6\textwidth]{fig12.2.png}
  \caption{Renormalization group flows near the free-field fixed point in
    scalar field theory for $d < 4$.}
  \label{fig:12.2}
\end{figure}

The diagram contributing to the correction $\Delta\lambda$ has the same
structure as the one-loop diagrams computed in Section~10.2.  In fact, these
are essentially the same diagrams, and differ only in whether the integrals
are carried out iteratively or all at once.  However, whereas the diagrams in
Section~10.2 had ultraviolet divergences, the corresponding diagram in
Wilson's approach is well defined and gives the coefficient of a simple
evolution equation of the coupling constant.  This transformation gives a
first example of the reinterpretation of ultraviolet divergences that we will
make in this chapter.

The transformation law \eqref{eq:12.28} implies that the renormalization
group flows near $\mathcal{L}_0$ have the form shown in Fig.~\ref{fig:12.1}(b),
with one slowly decaying direction.  If we follow the flows far enough, the
behavior should again be that of a free field.  This picture has the puzzling
implication that four-dimensional interacting $\phi^4$ theory does not exist
in the limit in which the cutoff goes to infinity.  We will discuss this
result further---and explain why it nevertheless makes sense to use $\phi^4$
theory as a model field theory---in Section~12.3.

Finally consider the case $d < 4$.  Now $\lambda$ becomes a relevant
parameter.  The renormalization group flows near $\mathcal{L}_0$ are shown in
Fig.~\ref{fig:12.2}.  The coupling constant grows as we go to lower momenta,
and the perturbation expansion breaks down at some scale.  This is the
situation in which the theory develops a nontrivial fixed point, as we will
see in Section~12.5.

\section{The Callan-Symanzik Equation}

How can we hope to obtain information on renormalization group flows from
the expressions for renormalized Green's functions, in which the cutoff has
already been taken to infinity?  We must first realize that renormalized
quantum field theories correspond to a restricted class of the full set of
possible Lagrangians that we considered in the previous section.  In Wilson's
language, a renormalized field theory with the cutoff taken arbitrarily large
corresponds to a trajectory that takes an arbitrarily long time to evolve to
a large value of the mass parameter.  Such a trajectory must, then, pass
arbitrarily close to a fixed point, which we will assume to be the
weak-coupling fixed point.  In the slow evolution past this fixed point, the
irrelevant operators in the original Lagrangian die away, and we are left
only with the relevant and marginal operators.  The coefficients of these
operators are in one-to-one correspondence with the parameters of the
renormalizable field theory.  Thus, in working with a renormalized field
theory, we are throwing away information on the evolution of irrelevant
perturbations, but keeping information on the flows of relevant and marginal
perturbations.

The flows of these parameters cannot be determined from the cutoff
dependence, because, in this framework, the cutoff has already been sent to
infinity.  However, we have an alternative, though more abstract, tool at our
disposal.  The parameters of a renormalized field theory are determined by a
set of renormalization conditions, which are applied at a certain momentum
scale (called the \emph{renormalization scale}).  By looking at how the
parameters of the theory depend on the renormalization scale, we can recover
the information contained in the renormalization group flows of the previous
section.

We consider first the specific case of $\phi^4$ theory in four dimensions,
where the coupling constant $\lambda$ is dimensionless and the corresponding
operator is marginal.  For simplicity, we will also assume that the mass term
$m^2$ has been adjusted to zero, so that the theory sits just at its critical
point.  We will perform this analysis in Minkowski space, using spacelike
reference momenta.  However, the analysis would be essentially identical if
carried out in Euclidean space.  If we wish to consider renormalization group
predictions at timelike momenta, we must consider the possibilities of new
singularities which make the analysis more complicated.  These include both
physical thresholds and the Sudakov double logarithms discussed in
Section~6.4.  We postpone discussion of these complications until
Chapters~17 and~18.

\subsection{Renormalization Conditions}

To define the theory properly, we must specify the renormalization
conditions.  In Chapter~10 we used a natural set of renormalization
conditions \eqref{eq:10.19} for $\phi^4$ theory, defined in terms of the
physical mass $m$.  However, in a theory where $m = 0$, these conditions
cannot be used because they lead to singularities in the counterterms.
(Consider, for example, the limit $m^2\to 0$ of Eq.~\eqref{eq:10.24}.)  To
avoid such singularities, we choose an arbitrary momentum scale $M$ and
impose the renormalization conditions at a spacelike momentum $p$ with
$p^2 = -M^2$:
\begin{equation}
  \begin{aligned}
    \text{(one-particle-irreducible two-point function)} &= 0
      \text{ at } p^2 = -M^2;\\
    \frac{d}{dp^2}\bigl( \text{two-point function} \bigr) &= 1
      \text{ at } p^2 = -M^2;\\
    \text{(four-point function)} &= -i\lambda
      \text{ at } (p_1+p_2)^2 = (p_1+p_3)^2 = (p_1+p_4)^2 = -M^2 .
  \end{aligned}
  \label{eq:12.30}
\end{equation}
The parameter $M$ is called the renormalization scale.  These conditions
define the values of the two- and four-point Green's functions at a certain
point and, in the process, remove all ultraviolet divergences.  Speaking
loosely, we say that we are ``defining the theory at the scale $M$.''  These
new renormalization conditions take some getting used to.  The second
condition, in particular, implies that the two-point Green's function has a
coefficient of 1 at the unphysical momentum $p^2 = -M^2$, rather than on
shell (at $p^2 = 0$):
\begin{equation}
  \langle\Omega|\phi(p)\phi(-p)|\Omega\rangle = \frac{i}{p^2}
  \qquad\text{at } p^2 = -M^2 .
\end{equation}
Here $\phi$ is the renormalized field, related to the bare field $\phi_0$ by a
scale factor that we again call $Z$:
\begin{equation}
  \phi = Z^{-1/2}\phi_0 .
  \label{eq:12.31}
\end{equation}
This $Z$, however, is not the residue of the physical pole in the two-point
Green's function of bare fields, as it was in Chapters~7 and~10.  Instead, we
now have
\begin{equation}
  \langle\Omega|\phi_0(p)\phi_0(-p)|\Omega\rangle = \frac{iZ}{p^2}
  \qquad\text{at } p^2 = -M^2 .
\end{equation}
The Feynman rules for renormalized perturbation theory are the same as in
Chapter~10, with the same relation between $Z$ and the counterterm
$\delta_Z$, $\delta_Z = Z - 1$.  Now, however, the counterterms $\delta_Z$ and
$\delta_\lambda$ must be adjusted to maintain the new conditions
\eqref{eq:12.30}.

The first renormalization condition in \eqref{eq:12.30} holds the physical
mass of the scalar field fixed at zero.  We saw in Chapter~10 that, in
$\phi^4$ theory, the one-loop propagator correction is momentum-independent
and is completely canceled by the mass renormalization counterterm.  At
two-loop order, however, the situation becomes more complicated, and the
propagator corrections require both mass and field strength
renormalizations.  In more general scalar field theories, such as the Yukawa
theory example considered at the end of Section~10.2, this complication
arises already at one-loop order.  Since the field strength renormalization
counterterm will play an important role in the discussion below, it will be
helpful to discuss briefly how we will treat this double subtraction.

The evaluation of propagator corrections has some special simplifications for
the case of a massless scalar field, which we consider here, and specifically
with the use of dimensional regularization.  Consider, for example, the
one-loop propagator correction in Yukawa theory.  In Section~10.2 we found an
expression of the form
\begin{equation}
  \Pi(p^2) \sim \frac{\Gamma(1-d/2)}{(4\pi)^{d/2}}\ \frac{p^2}{\Delta^{1-d/2}}
  ,
  \label{eq:12.32}
\end{equation}
where $\Delta$ is a linear combination of the fermion mass $m_f$ and $p^2$.
If we compute the diagram using massless propagators only, $\Delta$ is
proportional to $p^2$.  Expression \eqref{eq:12.32} has a pole at $d = 2$,
corresponding to the quadratically divergent mass renormalization.  However,
the residue of this pole is independent of $p^2$, so we can completely cancel
the pole with the mass counterterm $\delta_m$.  This allows us to
analytically continue \eqref{eq:12.32} to $d = 4$.  Then this expression
takes the form
\begin{equation}
  \Pi(p^2) \sim \frac{\Gamma(2-d/2)}{(4\pi)^{d/2}}\ \frac{p^2}{p^2}
  = \frac{p^2}{(4\pi)^2}\ \frac{2}{4-d} ,
  \label{eq:12.33}
\end{equation}
and gives no additional mass shift but only a field strength renormalization.
The remaining divergence is canceled by the counterterm $\delta_Z$.  If we
adopt the rule that we should simply continue expressions of the form
\eqref{eq:12.32} to $d = 4$, we can forget about the counterterm $\delta_m$
altogether.

In a regularization scheme with a momentum cutoff, the contributions to
$\delta_m$ and $\delta_Z$ become tangled up with one another.  Then it is more
awkward to define the massless limit.  In the following discussion, we will
assume the use of dimensional regularization.  However, to emphasize the
physical role of the cutoff, we will write expressions of the form
\eqref{eq:12.33} as
\begin{equation}
  \Pi(p^2) \sim p^2\ \log\Lambda^2 + (\text{finite terms}) .
  \label{eq:12.34}
\end{equation}
The logarithmically divergent terms proportional to $p^2$ will agree with the
divergences obtained with a momentum cutoff; the constant terms will not
agree, but these will drop out of our final results.

In $\phi^4$ theory, where the one-loop propagator correction is
momentum-independent, the one-loop diagram is simply set to zero by this
prescription.  Then the preceding analysis applies to the two-loop and higher
correction terms.

The generalization of the analysis of this section to massive scalar field
theory requires some additional formalism, which we postpone to
Section~12.5.

\subsection{The Callan-Symanzik Equation}

In the renormalization conditions \eqref{eq:12.30}, the renormalization scale
$M$ is arbitrary.  We could just as well have defined the same theory at a
different scale $M'$.  By ``the same theory'', we mean a theory whose bare
Green's functions,
\begin{equation}
  \langle\Omega|T\phi_0(x_1)\phi_0(x_2)\cdots\phi_0(x_n)|\Omega\rangle ,
\end{equation}
are given by the same functions of the bare coupling constant $\lambda_0$ and
the cutoff $\Lambda$.  These functions make no reference to $M$.  The
dependence on $M$ enters only when we remove the cutoff dependence by
rescaling the fields and eliminating $\lambda_0$ in favor of the renormalized
coupling $\lambda$.  The renormalized Green's functions are numerically equal
to the bare Green's functions, up to a rescaling by powers of the field
strength renormalization $Z$:
\begin{equation}
  \langle\Omega|T\phi(x_1)\phi(x_2)\cdots\phi(x_n)|\Omega\rangle
  = Z^{-n/2}\ \langle\Omega|T\phi_0(x_1)\phi_0(x_2)\cdots\phi_0(x_n)
    |\Omega\rangle .
  \label{eq:12.35}
\end{equation}
The renormalized Green's functions could be defined equally well at another
scale $M'$, using a new renormalized coupling $\lambda'$ and a new rescaling
factor $Z'$.

Let us write more explicitly the effect of an infinitesimal shift of $M$.
Let $G^{(n)}(x_1, \ldots, x_n)$ be the connected $n$-point function, computed
in renormalized perturbation theory:
\begin{equation}
  G^{(n)}(x_1, \ldots, x_n) = \langle\Omega|T\phi(x_1)\cdots\phi(x_n)
    |\Omega\rangle_{\mathrm{connected}} .
  \label{eq:12.36}
\end{equation}
Now suppose that we shift $M$ by $\delta M$.  There is a corresponding shift
in the coupling constant and the field strength such that the bare Green's
functions remain fixed:
\begin{equation}
  M \to M + \delta M,\qquad \lambda \to \lambda + \delta\lambda,\qquad
  \phi \to (1+\delta\eta)\phi .
  \label{eq:12.37}
\end{equation}
Then the shift in any renormalized Green's function is simply that induced by
the field rescaling,
\begin{equation}
  G^{(n)} \to (1+n\delta\eta)G^{(n)} .
\end{equation}
If we think of $G^{(n)}$ as a function of $M$ and $\lambda$, we can write this
transformation as
\begin{equation}
  \delta G^{(n)} = \Bigl( \delta M\frac{\partial}{\partial M}
    + \delta\lambda\frac{\partial}{\partial\lambda} \Bigr)G^{(n)}
  = n\delta\eta\,G^{(n)} .
  \label{eq:12.38}
\end{equation}
Rather than writing this relation in terms of $\delta\lambda$ and
$\delta\eta$, it is conventional to define the dimensionless parameters
\begin{equation}
  \beta \equiv \frac{M}{\delta M}\,\delta\lambda,\qquad
  \gamma \equiv \frac{M}{\delta M}\,\delta\eta .
  \label{eq:12.39}
\end{equation}
Making these substitutions in Eq.~\eqref{eq:12.38} and multiplying through by
$M/\delta M$, we obtain
\begin{equation}
  \Bigl[ M\frac{\partial}{\partial M} + \beta\frac{\partial}{\partial\lambda}
    + n\gamma \Bigr]\ G^{(n)}(x_i, \ldots) = 0 .
  \label{eq:12.40}
\end{equation}
The parameters $\beta$ and $\gamma$ are the same for every $n$, and must be
independent of the $x_i$.  Since the Green's function $G^{(n)}$ is
renormalized, $\beta$ and $\gamma$ cannot depend on the cutoff, and hence, by
dimensional analysis, these functions cannot depend on $M$.  Therefore they
are functions only of the dimensionless variable $\lambda$.  We conclude that
any Green's function of massless $\phi^4$ theory must satisfy
\begin{equation}
  \Bigl[ M\frac{\partial}{\partial M} + \beta(\lambda)\frac{\partial}
    {\partial\lambda} + n\gamma(\lambda) \Bigr]\ G^{(n)}\bigl( \{x_i\}; M,
      \lambda \bigr) = 0 .
  \label{eq:12.41}
\end{equation}
This relation is called the \emph{Callan-Symanzik equation}.%
\footnote{The equation was derived independently by C.~G.~Callan,
\emph{Phys.\ Rev.} \textbf{D2}, 1541 (1970) and K.~Symanzik,
\emph{Comm.\ Math.\ Phys.} \textbf{18}, 227 (1970).}  It asserts that there
exist two universal functions $\beta(\lambda)$ and $\gamma(\lambda)$, related
to the shifts in the coupling constant and field strength, that compensate
for the shift in the renormalization scale $M$.

The preceding argument generalizes without difficulty to other massless
theories with dimensionless couplings.  In theories with multiple fields and
couplings, there is a $\gamma$ term for each field and a $\beta$ term for
each coupling.  For example, we can define QED at zero electron mass by
introducing a renormalization scale as in Eqs.~\eqref{eq:12.30}.  The
renormalization conditions for the propagators are applied at $p^2 = -M^2$,
and those for the vertex at a point where all three invariants are of order
$-M^2$.  Then the renormalized Green's functions of this theory satisfy the
Callan-Symanzik equation
\begin{equation}
  \Bigl[ M\frac{\partial}{\partial M} + \beta(e)\frac{\partial}{\partial e}
    + n\gamma_2(e) + m\gamma_3(e) \Bigr]\ G^{(n,m)}\bigl( \{x_i\}; M, e \bigr)
  = 0 ,
  \label{eq:12.42}
\end{equation}
where $n$ and $m$ are, respectively, the number of electron and photon fields
in the Green's function $G^{(n,m)}$ and $\gamma_2$ and $\gamma_3$ are the
rescaling functions of the electron and photon fields.

\subsection{Computation of $\beta$ and $\gamma$}

Before we work out the implications of the Callan-Symanzik equation, let us
look more closely at the functions $\beta$ and $\gamma$ that appear in it.
From their definitions \eqref{eq:12.39}, we see that they are proportional to
the shift in the coupling constant and the shift in the field normalization,
respectively, when the renormalization scale $M$ is increased.  The behavior
of the coupling constant as a function of $M$ is of particular interest,
since it determines the strength of the interaction and the conditions under
which perturbation theory is valid.  We will see in the next section that the
shift in the field strength is also reflected directly in the values of
Green's functions.

The easiest way to compute the Callan-Symanzik functions is to begin with the
expressions for the counterterms.  For the $\gamma$ function, we consider the
renormalized two-point function.  At one-loop order, it has the general form
\begin{equation}
  G^{(2)}(p) = \frac{i}{p^2}\ \Bigl( 1 + \delta_Z + i\Pi(p^2) + \cdots
    \Bigr),
  \label{eq:12.49}
\end{equation}
where the one-loop diagram $\Pi(p^2)$ contains a logarithmically divergent
term proportional to $p^2$.  Applying the Callan-Symanzik equation
\eqref{eq:12.41} with $n = 2$ and neglecting the $\beta$ term (which is
always smaller by at least one power of the coupling constant), we find
\begin{equation}
  -M\frac{\partial}{\partial M}\,\delta_Z + 2\gamma = 0,
\end{equation}
or
\begin{equation}
  \gamma = \frac{1}{2}M\frac{\partial}{\partial M}\,\delta_Z
    \qquad \text{(to lowest order)} .
  \label{eq:12.50}
\end{equation}
To make this result more explicit, note that the counterterm must be
\begin{equation}
  \delta_Z = A\,\lambda^2\log\Lambda + \text{finite}
\end{equation}
in order to cancel the divergent logarithm in $G^{(2)}$.  Thus $\gamma$ is
simply the coefficient of the logarithm:
\begin{equation}
  \gamma = -A\,\lambda^2 \qquad \text{(to lowest order)} .
  \label{eq:12.51}
\end{equation}
In most theories (e.g., Yukawa theory or QED), the first logarithmic
divergence in $\delta_Z$ occurs at the one-loop level.  However, even in
$\phi^4$ theory, formulae \eqref{eq:12.50} and \eqref{eq:12.51} are true for
the first nonvanishing term in $\delta_Z$, in this case the two-loop
contribution.%
\footnote{At one loop, formula \eqref{eq:12.33} implies that we can also
identify $\lambda$ as the coefficient of $2/(4-d)$ in the 1PI self-energy, in
the limit $d\to 4$.  This relation changes in higher loops.  However,
Eq.~\eqref{eq:12.50} remains correct.}  By replacing the scalar field
propagator $(i/p^2)$ with a fermion propagator $(i\slashed{p}/p^2)$, we could
repeat this argument line for line to compute the $\gamma$ function for a
fermion field in terms of its field strength counterterm $\delta_{Z_2}$.

We can derive similar expressions for the $\beta$ function of a generic
dimensionless coupling constant $g$, associated with an $n$-point vertex.
Taking propagator corrections into account, the full connected Green's
function, to one-loop order, has the general form
\begin{equation}
  \begin{aligned}
    G^{(n)}(p_i) &= \Bigl( \prod_i \frac{i}{p_i^2} \Bigr)\ \Bigl[
      -i\delta g - i\Bigl( B\log\frac{\Lambda^2}{-p^2} - \delta_g \Bigr)
      + i\Bigl( \sum_i \delta_{Z_i}\log\frac{\Lambda^2}{p_i^2}
        - \delta_{Z_i} \Bigr) \Bigr]\\
    &= \Bigl( \prod_i \frac{i}{p_i^2} \Bigr)\ \Bigl[ -ig - i\Bigl(
      B - g\sum_i A_i \Bigr)\log\frac{\Lambda^2}{-p^2} + \text{finite terms}
      \Bigr] .
  \end{aligned}
  \label{eq:12.52}
\end{equation}
In this expression, $p_i$ are the momenta on the external legs, and $p^2$
represents a typical invariant built from these momenta.  We assume that
renormalization conditions are applied at a point where all such invariants
are spacelike and of order $-M^2$.  The $M$ dependence of this expression
comes from the counterterms $\delta_g$ and $\delta_{Z_i}$.  Applying the
Callan-Symanzik equation, we obtain
\begin{equation}
  \beta(g) = M\frac{\partial}{\partial M}\Bigl( -\delta g
    + g\sum_i \delta_{Z_i} \Bigr) \qquad \text{(to lowest order)} .
  \label{eq:12.53}
\end{equation}
To be more explicit, we note that
\begin{equation}
  \delta_g = -B\log\frac{\Lambda^2}{M^2} + \text{finite} .
\end{equation}
Thus the $\beta$ function is just a combination of the coefficients of the
divergent logarithms:
\begin{equation}
  \beta(g) = -2B - g\sum_i A_i \qquad \text{(to lowest order)} .
  \label{eq:12.54}
\end{equation}
Notice that the finite parts of counterterms are independent of $M$ and
therefore never contribute to $\beta$ or $\gamma$.  This means that, to
compute the leading terms in the Callan-Symanzik functions, we needn't be
too precise in specifying renormalization conditions: Any momentum scale of
order $M^2$ will yield the same results.  The divergent parts of the
counterterms can be estimated simply by setting all invariants inside of
logarithms equal to $M^2$, as we did above in our expression for the $n$-point
Green's function.

As in the computation of $\gamma$, this argument can be applied almost
without change to coupling constants for fields with spin.  In Yukawa theory,
for example, we consider the three-point function with one incoming fermion,
one outgoing fermion, and one scalar, with momenta $p_1$, $p_2$, and $p_3$,
respectively.  Then the tree-level expression for the three-point function
is
\begin{equation}
  \frac{i}{\slashed{p}_1}\ (-ig)\ \frac{i}{\slashed{p}_2}\ \frac{i}{p_3^2}
  .
  \label{eq:12.55}
\end{equation}
The one-loop corrections replace the quantity $(-ig)$ by the expression in
brackets in Eq.~\eqref{eq:12.52}.  Then formulae \eqref{eq:12.53} and
\eqref{eq:12.54} hold also for the $\beta$ function of this theory.

Similar expressions also apply in QED, though there are a number of small
complications.  The first comes in computing the $\gamma$ function for the
photon propagator.  In Eq.~(7.74), we saw that the general form of
the photon propagator in Feynman gauge is
\begin{equation}
  D^{\mu\nu}(q) = D(q)\Bigl( -g^{\mu\nu} + \frac{q^\mu q^\nu}{q^2} \Bigr)
    + \frac{1}{\xi}\frac{q^\mu q^\nu}{q^2} .
  \label{eq:12.56}
\end{equation}
The coefficient of the last term in \eqref{eq:12.56} depends on the gauge.
Fortunately, this term drops out of all gauge-invariant observables.  Thus it
makes sense to concentrate on the first term, projecting all external photons
onto their transverse components.  Projecting the photon propagator, we see
that $D(q)$ satisfies the Callan-Symanzik equation.  Since the corrections to
this function have the form \eqref{eq:12.49}, the arguments following that
formula are valid for photons as well as for electrons and scalars.  Thus, to
leading order,
\begin{equation}
  \gamma_3(e) = \frac{1}{2}M\frac{\partial}{\partial M}\,\delta_3 ,
  \label{eq:12.57}
\end{equation}
where $\delta_2$ and $\delta_3$ are the counterterms defined in
Section~10.3.

Similarly, we may consider the three-point connected Green's function
$\langle\psi(p_1)\bar{\psi}(p_2)A^\mu(q)\rangle$, projected onto transverse
components of the photon.  At leading order, this function equals
\begin{equation}
  \frac{i}{\slashed{p}_1}\ (-ie\gamma^\mu)\ \frac{i}{\slashed{p}_2}\
  \Bigl( -g^{\mu\nu} + \frac{q^\mu q^\nu}{q^2} \Bigr) .
\end{equation}
The divergent one-loop corrections have the same form, with $(-ie)$ replaced
by logarithmically divergent terms.  Thus, Eq.~\eqref{eq:12.53} gives the
lowest-order expression for the $\beta$ function:
\begin{equation}
  \beta(e) = M\frac{\partial}{\partial M}\Bigl( -e\delta_1 + e\delta_2
    + \frac{e}{2}\delta_3 \Bigr) .
  \label{eq:12.58}
\end{equation}
To find explicit expressions for the Callan-Symanzik functions of QED, we
must write expressions for the counterterms $\delta_1$, $\delta_2$, $\delta_3$.
In Section~10.3, we evaluated these counterterms using on-shell
renormalization conditions with massive fermions.  We must now re-evaluate
these terms for massless fermions and renormalization at $-M^2$.
Fortunately, we need only evaluate the logarithmically divergent pieces of
these counterterms, which are identical in the two cases.  Reading from
Eqs.~\eqref{eq:10.43} and \eqref{eq:10.44}, we find
\begin{equation}
  \delta_2 = \frac{-e^2}{16\pi^2}\ \frac{\Gamma(2-d/2)}{(M^2)^{2-d/2}},\qquad
  \delta_3 = \frac{-e^2}{48\pi^2}\ \frac{4\Gamma(2-d/2)}{(M^2)^{2-d/2}}
  \ \cdots .
  \label{eq:12.59}
\end{equation}
Using formulae \eqref{eq:12.57} and \eqref{eq:12.59}, we obtain at leading
order
\begin{equation}
  \gamma_2(e) = \frac{e^2}{16\pi^2},\qquad
  \gamma_3(e) = \frac{e^2}{48\pi^2} .
  \label{eq:12.60}
\end{equation}
And from Eq.~\eqref{eq:12.58}, we find
\begin{equation}
  \beta(e) = \frac{e^3}{12\pi^2} .
  \label{eq:12.61}
\end{equation}
It is important to remember that the expression we have used for $\delta_2$
explicitly assumes the use of Feynman gauge.  In fact, $\gamma_2$ depends on
the gauge parameter, and this makes sense, because Green's functions of
individual $\psi$ and $\bar{\psi}$ fields are not gauge invariant.  On the
other hand, the QED vacuum polarization, and therefore $\gamma_3$ and
$\beta$, are gauge invariant.

\subsection{The Meaning of $\beta$ and $\gamma$}

We can obtain a deeper insight into the nature of $\beta$ and $\gamma$ by
expressing them in terms of the parameters of bare perturbation theory: $Z$,
$\lambda_0$, and $\Lambda$ for the case of $\phi^4$ theory.

First recall that the bare and renormalized field are related by
\begin{equation}
  \phi(p) = Z^{1/2}(M)\,\phi_0(p) .
  \label{eq:12.62}
\end{equation}
This equation expresses the dependence of the field rescaling on $M$.  If $M$
is increased by $\delta M$, the renormalized field is shifted by
\begin{equation}
  \frac{Z(M+\delta M)^{-1/2}}{Z(M)^{-1/2}} - 1 .
\end{equation}
Hence our original definition \eqref{eq:12.39} of $\gamma$ gives us
immediately
\begin{equation}
  \gamma(\lambda) = \frac{1}{2}M\frac{\partial}{\partial M}\log Z .
  \label{eq:12.63}
\end{equation}
Since $\delta_Z = Z - 1$ (Eq.~\eqref{eq:10.17}), this formula is in agreement
with \eqref{eq:12.50} to leading order.  Formula \eqref{eq:12.63}, however, is
an exact relation.  This expression clarifies the relation of $\gamma$ to the
field strength rescaling.  However, it obscures the fact that $\gamma$ is
independent of the cutoff $\Lambda$.  To understand this aspect of $\gamma$,
we have to go back to the original definition of this function in terms of
renormalized Green's functions, whose cutoff independence follows from the
renormalizability of the theory.

Similarly, we can find an instructive expression for $\beta$ in terms of the
bare coupling constant.  The bare coupling constant is related to the
renormalized one by
\begin{equation}
  \lambda_0 = Z^{-2}\,\lambda\,M^{4-d}\,\cdots ,
  \label{eq:12.64}
\end{equation}
with the exact form depending on the renormalization scheme.  Since
$\lambda_0$ is independent of $M$, we can differentiate \eqref{eq:12.64}
with respect to $M$ to obtain a relation that determines $\beta(\lambda)$.
The result is
\begin{equation}
  \beta(\lambda) = M\frac{\partial\lambda}{\partial M}
  = -\lambda\,\Bigl( \frac{d-4}{2} \Bigr) + \cdots ,
  \label{eq:12.66}
\end{equation}
up to corrections from the $Z$ factors.  In four dimensions, this gives the
explicit one-loop result
\begin{equation}
  \beta(\lambda) = \frac{3\lambda^2}{16\pi^2} .
  \label{eq:12.46}
\end{equation}

\section{Evolution of Coupling Constants}

We are now in a position to solve the Callan-Symanzik equation and extract
the physical content of the renormalization group.  Let us consider the
solution of the Callan-Symanzik equation for a family of Green's functions
related by a uniform rescaling of the momenta.  For the two-point function of
$\phi^4$ theory, with $p^2 = -P^2$ spacelike, the Callan-Symanzik equation
takes the form
\begin{equation}
  \Bigl[ P\frac{\partial}{\partial P} - \beta(\lambda)\frac{\partial}
    {\partial\lambda} + 2\gamma(\lambda) \Bigr]\ G^{(2)}(P; \lambda) = 0 ,
  \label{eq:12.70}
\end{equation}
where we have used the fact that $G^{(2)}$ depends only on the ratio
$P/M$ and $\lambda$.  The solution of this equation can be written as
\begin{equation}
  G^{(2)}(P; \lambda) = \bar{G}^{(2)}\bigl( \bar{\lambda}(P) \bigr)\
    \exp\Bigl[ -2\int_{M}^{P} \frac{dp'}{p'}\ \gamma\bigl(
      \bar{\lambda}(p') \bigr) \Bigr] ,
  \label{eq:12.76}
\end{equation}
where the \emph{running coupling constant} $\bar{\lambda}(p)$ is defined as
the solution of the differential equation
\begin{equation}
  \frac{d}{d\log(p/M)}\ \bar{\lambda}(p) = \beta\bigl( \bar{\lambda}(p)
    \bigr),\qquad \bar{\lambda}(M) = \lambda ,
  \label{eq:12.73}
\end{equation}
and $\bar{G}^{(2)}$ is determined by matching the solution at $p = M$:
\begin{equation}
  \bar{G}^{(2)}(\lambda) = G^{(2)}(M; \lambda) .
\end{equation}
For the two-point function in $\phi^4$ theory, this matching is rather
trivial: $G^{(2)} = 1 + \mathcal{O}(\lambda)$.

The preceding analysis can be applied to any family of Green's functions that
are related by uniform rescaling of the momenta.  Consider, for example, the
connected four-point function of $\phi^4$ theory evaluated at spacelike
momenta $p_i$ such that $p_i^2 = -P^2$, $p_i\cdot p_j = 0$, so that $s$, $t$,
and $u$ are of order $-P^2$.  To leading order in perturbation theory, this
function is given by
\begin{equation}
  G^{(4)}(P) = \Bigl( \frac{i}{P^2} \Bigr)^4\ \hat{G}^{(4)}(\lambda) .
  \label{eq:12.77}
\end{equation}
Using the fact that $G^{(4)}$ has dimensions of $(\mathrm{mass})^{-8}$, we can
exchange $M$ for $P$ in the Callan-Symanzik equation and write this equation
as
\begin{equation}
  \Bigl[ P\frac{\partial}{\partial P} - \beta(\lambda)\frac{\partial}
    {\partial\lambda} + 8 - 4\gamma(\lambda) \Bigr]\ \hat{G}^{(4)}(P; \lambda)
  = 0 .
  \label{eq:12.78}
\end{equation}
The solution to this equation is
\begin{equation}
  \hat{G}^{(4)}(P; \lambda) = \bar{G}^{(4)}\bigl( \bar{\lambda}(P) \bigr)\
    \exp\Bigl[ 4\int_{M}^{P} \frac{dp'}{p'}\ \gamma\bigl(
      \bar{\lambda}(p') \bigr) \Bigr] .
  \label{eq:12.79}
\end{equation}
This formula must agree with \eqref{eq:12.77} to leading order in $\lambda$;
this matching requires that
\begin{equation}
  \bar{G}^{(4)}\bigl( \bar{\lambda}(P) \bigr) = -i\bar{\lambda}
    + \mathcal{O}(\bar{\lambda}^2) .
  \label{eq:12.80}
\end{equation}

We can now see the physical implication of the Callan-Symanzik equation.
The ordinary Feynman perturbation series for a Green's function depends both
on the coupling constant $\lambda$ and on the dimensionless parameter
$\log(-p^2/M^2)$.  The perturbation theory can be badly behaved even when
$\lambda$ is small if the ratio $p^2/M^2$ is large.  The solutions
\eqref{eq:12.76} and \eqref{eq:12.79} reorganize this dependence into a
function of the running coupling constant and an exponential scale factor.
We consider these two pieces in turn.

The first factor in Eqs.~\eqref{eq:12.76} and \eqref{eq:12.79} is a function
of the running coupling constant, evaluated at the momentum scale $p$.  If
$p$ were of order $M$, the renormalization scale, this function would
essentially be the ordinary perturbative evaluation of the Green's function.
The results \eqref{eq:12.76} and \eqref{eq:12.79} instruct us to make use of
this same expression at the scale $p$, but to replace $\lambda$ with a new
coupling constant $\bar{\lambda}$ appropriate to that scale.  Thus, the
running coupling constant $\bar{\lambda}(p)$ is precisely the effective
coupling constant of the renormalization group flow.  This interpretation is
particularly clear in the solution \eqref{eq:12.79} for $\hat{G}^{(4)}(P)$,
since this function directly measures the strength of the $\phi^4$ coupling
constant.

The exponential factor in Eqs.~\eqref{eq:12.76} and \eqref{eq:12.79} has an
equally simple interpretation: It is the accumulated field strength rescaling
of the correlation function from the reference point $M$ to the actual
momentum $p$ at which the Green's function is evaluated.  This factor
receives a multiplicative contribution from each intermediate scale between
$M$ and $p$.  Each of these contributions is, appropriately, computed using
the running coupling constant at that particular scale.

As a check on these formal arguments, we can use the explicit form of the
$\beta$ function of $\phi^4$ theory found in Eq.~\eqref{eq:12.46} and the
renormalization group equation \eqref{eq:12.73} to evaluate the running
coupling constant of $\phi^4$ theory.  This running coupling constant
satisfies the differential equation
\begin{equation}
  \frac{d}{d\log(p/M)}\,\bar{\lambda} = \frac{3\bar{\lambda}^2}{16\pi^2},
  \qquad\text{with } \bar{\lambda}(M;\lambda) = \lambda .
  \label{eq:12.81}
\end{equation}
Integrating, we find
\begin{equation}
  \frac{1}{\bar{\lambda}} - \frac{1}{\lambda} = -\frac{3}{16\pi^2}\log
    \frac{p}{M},
\end{equation}
and thus,
\begin{equation}
  \bar{\lambda}(p) = \frac{\lambda}{1 - (3\lambda/16\pi^2)\log(p/M)} .
  \label{eq:12.82}
\end{equation}
Many properties of the solution to the Callan-Symanzik equation are visible
in this relation.  First, the expansion of this formula for $\bar{\lambda}$ to
order $\lambda^2$ agrees precisely with Eq.~\eqref{eq:12.28}, the rate of the
renormalization group flow from Wilson's method.  Second, this expression for
the running coupling constant goes to zero at a logarithmic rate as $p\to 0$.
This coincides with our expectation that a positive value for the $\beta$
function should imply an effective coupling that becomes stronger at large
momenta and weaker at small momenta.

If we expand the running coupling constant $\bar{\lambda}(p)$ in powers of
$\lambda$, we find that the successive powers of the coupling constant are
multiplied by powers of logarithms,
\begin{equation}
  \lambda^{n+1}\bigl( \log p/M \bigr)^n ,
\end{equation}
which become large and invalidate a simple perturbation expansion for $p$
much greater or much less than $M$.  We have seen this problem of large
logarithms arising several times in our diagram calculations, and we have
remarked on it specifically as a problem in the discussion following
Eq.~\eqref{eq:11.81}.  We now see that the renormalization group gives a
partial solution to this problem.  In this example, and in many others that
we will study, the Callan-Symanzik equation tells us how to sum these large
logarithms into the running coupling constant and multiplicative rescalings.
If the running coupling constant becomes large, as happens in $\phi^4$ theory
for $p\to\infty$, the perturbation expansion will break down anyway, and we
will need more advanced methods.  However, if the running coupling constant
becomes small, as for $\phi^4$ theory as $p\to 0$, we will have successfully
organized the powers of logarithms into a meaningful and controlled
expression.  The specific problem posed at the end of Section~11.4 will be
solved explicitly by this method in Section~13.2.

\subsection{An Application to QED}

For a more concrete application of the Callan-Symanzik equation, we can look
again at the electromagnetic potential between static charges, $V(\mathbf{x})$,
which we studied in Section~7.5.  At very short distances or at large
momenta, we can ignore the electron mass in the computation of QED
corrections to this potential.  In this approximation, the potential should
obey the Callan-Symanzik equation of massless QED.  We could write this
equation either for $V(\mathbf{x})$ itself or for its Fourier transform; we
choose to work in Fourier space in order to make contact more easily with the
results of Section~7.5.

We define the massless limit of QED by specifying a renormalization scale $M$
at which the renormalized coupling $e_r$ is defined.  If $M$ is taken close to
the electron mass $m$, at the point where the massless approximation is just
becoming valid, then the value of $e_r$ will be close to the physical electron
charge $e$.  The potential between static charges is a measurable energy, so
its normalization is unambiguous and is not shifted from one renormalization
point to another.  Thus the Callan-Symanzik equation for the Fourier
transform of the potential has no $\gamma$ term, being simply
\begin{equation}
  \Bigl[ M\frac{\partial}{\partial M} + \beta(e_r)\frac{\partial}
    {\partial e_r} \Bigr]\ V(q; M, e_r) = 0 .
  \label{eq:12.83}
\end{equation}
The Fourier transform of the potential has dimensions of $(\mathrm{mass})^{-2}$,
so we can trade dependence on $M$ for dependence on $q$ as in the scalar
field theory discussion above.  This gives
\begin{equation}
  \Bigl[ q\frac{\partial}{\partial q} - \beta(e_r)\frac{\partial}{\partial
    e_r} + 2 \Bigr]\ V(q; M, e_r) = 0 .
  \label{eq:12.84}
\end{equation}
Equation \eqref{eq:12.84} is almost the same as Eq.~\eqref{eq:12.66}, so we
can immediately write down the solution as a special case of
\eqref{eq:12.76}:
\begin{equation}
  V(q, e_r) = \frac{1}{q^2}\ \bar{V}\bigl( \bar{e}(q; e_r) \bigr) ,
  \label{eq:12.85}
\end{equation}
where $\bar{e}(q)$ is the solution of the renormalization group equation
\begin{equation}
  \frac{d}{d\log(q/M)}\ \bar{e} = \beta(\bar{e}),\qquad
  \bar{e}(M; e_r) = e_r .
  \label{eq:12.86}
\end{equation}
By comparing this formula for $V(q)$ to the leading-order result
$V(q) = e_r^2/q^2$, we can identify $\bar{V}(\bar{e}) = \bar{e}^2 +
\mathcal{O}(\bar{e}^4)$.  Then
\begin{equation}
  V(q, e_r) = \frac{\bar{e}^2(q)}{q^2} ,
  \label{eq:12.87}
\end{equation}
up to corrections that are suppressed by powers of $e^2$ and contain no
compensatory large logarithms of $q/M$.

To turn Eq.~\eqref{eq:12.87} into a completely explicit formula, we need only
solve the renormalization group equation \eqref{eq:12.86}.  Using the QED
$\beta$ function \eqref{eq:12.61}, we can integrate \eqref{eq:12.86} to find
\begin{equation}
  \frac{12\pi^2}{e^2}\Bigl( \frac{1}{\bar{e}^2} - \frac{1}{e_r^2} \Bigr)
  = \log\frac{q}{M} .
\end{equation}
This simplifies to
\begin{equation}
  \bar{e}^2(q) = \frac{e_r^2}{1 - (e_r^2/6\pi^2)\log(q/M)} .
  \label{eq:12.88}
\end{equation}
This result is almost identical to the formula for the effective electric
charge that we found in Eq.~\eqref{eq:7.96}.  To cement the identification,
set $M$ to be of order the electron mass, $M^2 = \Lambda m^2$, and
approximate $e_r$ at this point by $e$, with $\alpha = e^2/4\pi$.  Then
Eq.~\eqref{eq:12.88} takes the form
\begin{equation}
  \bar{\alpha}(q) = \frac{\alpha}{1 - (\alpha/3\pi)\log(-q^2/\Lambda m^2)} .
  \label{eq:12.89}
\end{equation}
The particular choice $\Lambda = \exp(5/3)$ reproduces Eq.~\eqref{eq:7.96}.
Of course, we could not find this exact correspondence without the detailed
one-loop calculation of Section~7.5.  Nevertheless, our present analysis
produces the correct asymptotic formula for the effective charge.
Furthermore, our present formalism can be applied to any renormalizable
quantum field theory; it does not rely on the special symmetries of QED that
we exploited in Section~7.5.

\subsection{Alternatives for the Running of Coupling Constants}

Now that we have computed the behavior of the running coupling constant in
two specific quantum field theories, let us consider more generally what
behaviors of the running coupling constant are possible in principle.  We
continue to restrict our discussion to renormalizable theories in the massless
limit, with a single dimensionless coupling constant $\lambda$.

By the arguments of the previous section, the Green's functions in any such
theory obey a Callan-Symanzik equation.  The solution of this equation
depends on a running coupling constant, $\bar{\lambda}(p)$, which satisfies a
differential equation
\begin{equation}
  \frac{d}{d\log(p/M)}\ \bar{\lambda} = \beta(\bar{\lambda}),
  \label{eq:12.90}
\end{equation}
in which the function $\beta(\lambda)$ is computable as a power series in the
coupling constant.  In the examples we have just discussed, the leading
coefficient in this power series was positive.  However, as a matter of
principle, three behaviors are possible in the region of small $\lambda$:
\begin{equation}
  \begin{aligned}
    (1)\quad &\beta(\lambda) > 0;\\
    (2)\quad &\beta(\lambda) = 0;\\
    (3)\quad &\beta(\lambda) < 0 .
  \end{aligned}
\end{equation}

In case (1), the coupling constant grows as the momentum increases, and the
theory becomes strongly coupled at large momenta.  This is the situation in
$\phi^4$ theory and QED.  The theory is \emph{asymptotically free} only in
case (3), where the coupling decreases at large momenta.  We will see in
Chapter~16 that non-Abelian gauge theories (such as QCD) fall into case (3).
Case (2) is the marginal situation, which requires the computation of higher
order terms in $\beta$ to determine the flow.

The behavior of the running coupling constant determines the validity of
perturbation theory at large momenta.  In case (1), perturbation theory
breaks down at some large momentum scale, and the theory must be defined by
additional input (such as the Landau pole in QED).  In case (3), the theory
becomes weakly coupled at large momenta, so that perturbative predictions are
reliable there.

If the $\beta$ function has a zero at a nonzero value $\lambda^{*}$ of the
coupling constant, we say that the theory has a \emph{nontrivial fixed point}
at $\lambda^{*}$.  Near such a fixed point, the running coupling approaches
$\lambda^{*}$ at a rate determined by the slope of $\beta$ at the fixed
point.  The behavior of Green's functions near the fixed point is governed by
the value of the anomalous dimension at the fixed point.  Such fixed points
play a central role in the theory of critical phenomena, as we will discuss
in Chapter~13.  The physical predictions (such as the critical exponents)
obtained near the fixed point---the value of $\lambda^{*}$ itself, the slope
$B$ at the zero, and the value of the anomalous dimension at the fixed
point---should all be independent of the conventions used to compute $\beta$
and $\gamma$.

\section{Renormalization of Local Operators}

The analysis of the previous two sections has been restricted to quantum
field theories with only dimensionless coefficients, that is, strictly
renormalizable field theories in the massless limit.  It is not difficult to
generalize this formalism to theories with mass terms and other operators
whose coefficients have mass dimension.  However, it is worthwhile to first
devote some attention to an intermediate step, by analyzing the renormalization
group properties of matrix elements of local operators.  This is an
interesting problem in its own right, and we will devote considerable space
to the applications of this formalism in Chapter~18.

Matrix elements of local operators appear often in quantum field theory
calculations.  Typically one considers a set of interacting particles that
couple weakly to an additional particle, which mediates new forces.
Consider, for example, the theory of strongly interacting quarks perturbed by
the effects of weak decay processes.  The weak interaction is mediated by a
massive vector boson, the $W$.  Let us write the interaction of the quarks
with the $W$ very schematically as
\begin{equation}
  \delta\mathcal{L} = \frac{g}{\sqrt{2}}\ \bar{\psi}\gamma^\mu(1-\gamma^5)\psi\,
    W_\mu ,
  \label{eq:12.98}
\end{equation}
and assign the $W$ boson the propagator
\begin{equation}
  \frac{-i}{q^2-m_W^2}\ \Bigl( g^{\mu\nu} - \frac{q^\mu q^\nu}{m_W^2}
    \Bigr) .
\end{equation}
(We will discuss this interaction more correctly in Section~18.2 and in
Chapter~20.)  Exchange of a $W$ boson leads to the interaction shown in
Fig.~\ref{fig:12.5}.  For momentum transfers small compared to $m_W$, we can
ignore the $q^2$ in the $W$ propagator and write this interaction as the
matrix element of the operator
\begin{equation}
  -\frac{g^2}{2m_W^2}\ \mathcal{O}(x) ,
  \label{eq:12.100}
\end{equation}
where
\begin{equation}
  \mathcal{O}(x) = \bar{\psi}\gamma^\mu(1-\gamma^5)\psi\,
    \bar{\psi}\gamma_\mu(1-\gamma^5)\psi .
\end{equation}
In the spirit of Wilson's renormalization group procedure, we can say that,
on distance scales larger than $m_W^{-1}$, the $W$ boson can be integrated
out, leaving over the interaction \eqref{eq:12.100}.

\begin{figure}[t]
  \centering
  \includegraphics[width=0.5\textwidth]{fig12.5.png}
  \caption{Interaction of quarks generated by the exchange of a $W$ boson.}
  \label{fig:12.5}
\end{figure}

How would we analyze the effects of the operator \eqref{eq:12.100} on
strongly interacting particles composed of quarks and antiquarks?  A useful
way to begin is to compute the Green's function of the operator
$\mathcal{O}$ together with fields that create and destroy quarks.  If we
approximate the theory of quarks by a theory of free fermions, it is easy to
compute these Green's functions; for example:
\begin{equation}
  \langle\psi(p_1)\bar{\psi}(-p_2)\psi(p_3)\bar{\psi}(-p_4)
    \mathcal{O}(0)\rangle
  = S_F(p_1)\gamma^\mu(1-\gamma^5)S_F(p_2)\ S_F(p_3)\gamma_\mu(1-\gamma^5)
    S_F(p_4) .
  \label{eq:12.101}
\end{equation}
However, in an interacting field theory, the answer will be much more
complicated.  Some of these complications will involve the low-energy
interactions of quarks, and we will leave them outside of the present
discussion.  However, in a renormalizable theory of quark interactions, one
will also find that Green's functions containing $\mathcal{O}$ have new
ultraviolet divergences.  The one-loop corrections to \eqref{eq:12.101} will
contain diagrams that evaluate to the right-hand side of \eqref{eq:12.101}
times a divergent integral.  These diagrams can be interpreted as field
strength renormalizations of the operator $\mathcal{O}$.  As with correlation
functions of elementary fields, we can obtain finite and well-defined matrix
elements of local operators only if we establish conventions for the
normalization of local operators and introduce operator rescalings in the
form of counterterms, order by order in perturbation theory, to preserve
these conventions.  More specifically, in a massless, renormalizable field
theory of the fermions $\psi$, we should make the convention that
Eq.~\eqref{eq:12.101} is exact at some spacelike normalization point for
which $p_1^2 = p_2^2 = p_3^2 = p_4^2 = -M^2$.  Then we should add a
counterterm of the form $\delta_{\mathcal{O}}\mathcal{O}(x)$, and adjust this
counterterm at each order of perturbation theory to insure that these
relations are preserved.  We refer to the operator satisfying the
normalization condition \eqref{eq:12.101} at $M^2$ as $\mathcal{O}_M$.

The renormalized operator $\mathcal{O}_M$ is a rescaled version of the
operator $\mathcal{O}_0$ built of bare fields,
\begin{equation}
  \mathcal{O}_0(x) = \bar{\psi}_0\gamma^\mu(1-\gamma^5)\psi_0\,
    \bar{\psi}_0\gamma_\mu(1-\gamma^5)\psi_0 .
  \label{eq:12.102}
\end{equation}
As we did for the elementary fields, we can write this relation as
\begin{equation}
  \mathcal{O}_0 = Z_{\mathcal{O}}(M)\,\mathcal{O}_M .
  \label{eq:12.103}
\end{equation}
This allows us to write the generalization of the relation \eqref{eq:12.35}
between Green's functions of bare and renormalized fields.  Let us return to
the language of scalar field theories and consider $\mathcal{O}(x)$ to be a
local operator in a scalar field theory.  Define
\begin{equation}
  G^{(n;1)}(p_1, \ldots, p_n; k) = \bigl\langle \phi(p_1)\cdots\phi(p_n)\,
    \mathcal{O}_M(k) \bigr\rangle .
  \label{eq:12.104}
\end{equation}
Then $G^{(n;1)}$ is related to a Green's function of bare fields by
\begin{equation}
  G^{(n;1)}(p_1, \ldots, p_n; k) = Z^{n/2}Z_{\mathcal{O}}(M)^{-1}\
    \bigl\langle \phi_0(p_1)\cdots\phi_0(p_n)\,\mathcal{O}_0(k) \bigr\rangle .
  \label{eq:12.105}
\end{equation}
Repeating the derivation of Eqs.~\eqref{eq:12.63} and \eqref{eq:12.64}, we
find that the Green's functions containing a local operator obey the
Callan-Symanzik equation
\begin{equation}
  \Bigl[ M\frac{\partial}{\partial M} + \beta(\lambda)\frac{\partial}
    {\partial\lambda} + n\gamma(\lambda) + \gamma_{\mathcal{O}}(\lambda) \Bigr]
    G^{(n;1)} = 0 ,
  \label{eq:12.106}
\end{equation}
where
\begin{equation}
  \gamma_{\mathcal{O}} = M\frac{\partial}{\partial M}\log Z_{\mathcal{O}}(M)
  .
  \label{eq:12.107}
\end{equation}

It often happens that a quantum field theory contains several operators with
the same quantum numbers.  For example, in quantum electrodynamics, the
operators $\bar{\psi}[\gamma^\mu D^\nu + \gamma^\nu D^\mu]\psi$ and
$F^{\mu\lambda}F^{\nu}_{\ \lambda}$ are both symmetric tensors with zero
electric charge; in addition, both operators have mass dimension 4.  Such
operators, with the same quantum numbers and the same mass dimension, can be
\emph{mixed} by quantum corrections.  For such a set of operators
$\{\mathcal{O}_i\}$, the relation of renormalized and bare operators must be
generalized to
\begin{equation}
  \mathcal{O}_i = Z^{ij}_{\mathcal{O}}(M)\,\mathcal{O}_j .
  \label{eq:12.108}
\end{equation}
This relation in turn implies that the anomalous dimension function
$\gamma_{\mathcal{O}}$ in the Callan-Symanzik equation must be generalized to
a matrix,
\begin{equation}
  \gamma_{\mathcal{O}} = \bigl[ Z_{\mathcal{O}}(M) \bigr]^{-1}\ M\frac{\partial}
    {\partial M}\ Z_{\mathcal{O}}(M) .
  \label{eq:12.109}
\end{equation}
Most of our applications of \eqref{eq:12.106} in Chapter~18 will require this
generalization.%
\footnote{Our assumption that we are working in a massless field theory
constrains the possibilities for operator mixing.  In a massive field theory,
operators of a given dimension can also mix with operators of lower
dimension.}

On the other hand, there are some operators for which the rescaling and
anomalous dimensions are especially simple.  If $\mathcal{O}$ is the quark
number current $\bar{\psi}\gamma^\mu\psi$, its normalization is fixed once and
for all because the associated charge
\begin{equation}
  Q = \int d^3x\ \bar{\psi}\gamma^0\psi
\end{equation}
is just the conserved integer number of quarks minus antiquarks in a given
state.  More generally, for any conserved current $J^\mu$,
$Z_J(M) = 1$ and $\gamma_J = 0$.  The same argument applies to the
energy-momentum tensor.  Thus, in the QED example above, the specific linear
combination
\begin{equation}
  T^{\mu\nu} = \bar{\psi}\bigl[ \gamma^\mu D^\nu + \gamma^\nu D^\mu \bigr]\psi
    - \cdots ,
  \label{eq:12.110}
\end{equation}
receives no rescaling and no anomalous dimension.  This linear combination of
operators must be an eigenvector of the matrix $\gamma^{ij}$ with eigenvalue
zero.

So far, our discussion of operator matrix elements has been rather abstract.
To make it more concrete, we will construct a formula for computing
$\gamma_{\mathcal{O}}$ to leading order from one-loop counterterms, and then
apply this formula to a simple example in $\phi^4$ theory.

To find a simple formula for $\gamma_{\mathcal{O}}$, we follow the same path
that took us from Eq.~\eqref{eq:12.52} to the formula \eqref{eq:12.53} for
the $\beta$ function.  Consider an operator whose normalization condition is
based on a Green's function with $m$ scalar fields:
\begin{equation}
  G^{(m;1)} = \bigl\langle \phi(p_1)\cdots\phi(p_m)\,\mathcal{O}_M(k)
    \bigr\rangle .
  \label{eq:12.111}
\end{equation}
To compute this Green's function to one-loop order, we find the set of
diagrams:
\begin{equation}
  \text{(tree level)} + \text{(one-loop vertex correction)} + \text{(external
    leg corrections)} + \text{(counterterm)} .
\end{equation}
The last diagram is the counterterm $\delta_{\mathcal{O}}$ needed to maintain
the renormalization condition.  Notice that the counterterm $\delta_Z$ also
appears.  If we insist that this sum of diagrams satisfies the Callan-Symanzik
equation \eqref{eq:12.106} to leading order in $\lambda$, we find,
analogously to \eqref{eq:12.53}, the relation
\begin{equation}
  \gamma_{\mathcal{O}}(\lambda) = M\frac{\partial}{\partial M}\Bigl(
    \delta_{\mathcal{O}} + \frac{m}{2}\delta_Z \Bigr) .
  \label{eq:12.112}
\end{equation}

As a specific example of the use of this formula, let us compute the
anomalous dimension $\gamma_{\mathcal{O}}$ of the mass operator $\phi^2$ in
$\phi^4$ theory.  There is a small subtlety involved in this computation.
The Feynman diagrams of $\phi^4$ theory generate an additive mass
renormalization, which must be removed by the mass counterterm at each order
in perturbation theory.  We would like to define the mass operator as a
perturbation which we can add to the massless theory defined in this way.  To
clarify the distinction between the underlying mass, which is renormalized to
zero, and the explicit mass perturbation, we will analyze a Green's function
of $\phi^2$ in which this operator carries a specific nonzero momentum.  We
thus choose to define the normalization of $\phi^2$ by the convention
\begin{equation}
  \bigl\langle \phi(p)\phi(q)\phi^2(k) \bigr\rangle
  = \frac{i}{p^2}\frac{i}{q^2}\frac{i}{k^2}\cdot 2
  \qquad\text{at } p^2 = q^2 = k^2 = -M^2 .
  \label{eq:12.113}
\end{equation}
The one-particle-irreducible one-loop correction to \eqref{eq:12.113} is
\begin{equation}
  \frac{i}{p^2}\frac{i}{q^2}\ \frac{i}{2}(-\lambda)\ \int
    \frac{d^dr}{(2\pi)^d}\ \frac{i}{r^2}\ \frac{i}{(k+r)^2}
  = \frac{i}{p^2}\frac{i}{q^2}\ \frac{\lambda}{2}\ \frac{\Gamma(2-d/2)}
    {(4\pi)^{d/2}}\ \frac{1}{(M^2)^{2-d/2}}\ \cdots ,
  \label{eq:12.114}
\end{equation}
where $\Delta$ is a function of the external momenta.  At $-M^2$, this
contribution must be canceled by a counterterm diagram,
\begin{equation}
  \frac{i}{p^2}\frac{i}{q^2}\ \delta_{\mathcal{O}}\ \frac{i}{k^2}\cdot 2 .
  \label{eq:12.115}
\end{equation}
Thus, the counterterm must be
\begin{equation}
  \delta_{\mathcal{O}} = -\frac{\lambda}{2}\ \frac{\Gamma(2-d/2)}
    {(4\pi)^2}\ \frac{1}{(M^2)^{2-d/2}} .
  \label{eq:12.116}
\end{equation}
Since $\delta_Z$ is finite to order $\lambda$, this is the only contribution
to \eqref{eq:12.112}, and we find
\begin{equation}
  \gamma_{\phi^2} = \frac{\lambda}{16\pi^2} .
  \label{eq:12.117}
\end{equation}

\section{Evolution of Mass Parameters}

We now discuss the renormalization group evolution of the mass parameter of a
scalar field theory.  In the preceding sections we have considered the
massless theory, in which the mass is adjusted to zero at the critical point.
We now generalize the analysis to include a mass term, which is a relevant
operator.  The formalism follows the same lines as before, but now the
Callan-Symanzik equation contains an additional term from the mass.

We consider the massive $\phi^4$ theory, defined by the renormalization
conditions \eqref{eq:12.30} generalized to include a nonzero mass.  The
renormalized Lagrangian contains a mass term $\tfrac{1}{2}m^2\phi^2$, and the
renormalization conditions are imposed at spacelike momenta of order $-M^2$.
The bare mass $m_0$ is related to the renormalized mass $m$ by
\begin{equation}
  m_0^2 = Z_m\,m^2 ,
\end{equation}
where $Z_m$ is a new renormalization constant.  The Callan-Symanzik equation
for the massive theory takes the form
\begin{equation}
  \Bigl[ M\frac{\partial}{\partial M} + \beta(\lambda)\frac{\partial}
    {\partial\lambda} - \gamma_m(\lambda)\,m^2\frac{\partial}{\partial m^2}
    + n\gamma(\lambda) \Bigr]\ G^{(n)} = 0 ,
  \label{eq:12.126}
\end{equation}
where $\gamma_m$ is the anomalous dimension of the mass operator, which is
related to $Z_m$ by
\begin{equation}
  \gamma_m = M\frac{\partial}{\partial M}\log Z_m .
\end{equation}
To leading order, we saw in the previous section that the mass operator
$\phi^2$ has anomalous dimension
\begin{equation}
  \gamma_m = \frac{\lambda}{16\pi^2} + \cdots ,
  \label{eq:12.127}
\end{equation}
for $\phi^4$ theory.

It is convenient to parametrize the mass by the dimensionless quantity
\begin{equation}
  \rho_m = \frac{m^2}{M^2} .
\end{equation}
Then the Callan-Symanzik equation \eqref{eq:12.126} can be written in terms
of $\rho_m$, and the running mass satisfies the equation
\begin{equation}
  \frac{d}{d\log(p/M)}\ \rho_m = \bigl[ -2 + \gamma_m(\lambda) \bigr]\rho_m .
  \label{eq:12.131}
\end{equation}
For a general field theory with several couplings $\lambda_i$ and masses
$m_i$, the evolution equations take the form
\begin{equation}
  \beta = (d-4)\lambda + \beta^{(4)}(\lambda) + \cdots,\qquad
  \beta_m = \bigl[ -2 + \gamma_m^{(4)} \bigr]m^2 + \cdots ,
  \label{eq:12.132}
\end{equation}
where the functions with a superscript $(4)$ are the four-dimensional results
obtained earlier in this section.

\subsection{Critical Exponents: A First Look}

As an application of the formalism of this section, let us calculate the
renormalization group flow of the coefficient of the mass operator in
$\phi^4$ theory.  This is found by integrating Eq.~\eqref{eq:12.126}, using
the value of $\beta_m$ from \eqref{eq:12.131}:
\begin{equation}
  \frac{d}{d\log(p/M)}\ \rho_m = \bigl[ -2 + \gamma_{\phi^2}(\lambda) \bigr]
    \rho_m .
  \label{eq:12.134}
\end{equation}
For $\lambda = 0$, this equation gives the trivial relation
\begin{equation}
  \rho_m(p) = \rho_m(M)\ \Bigl( \frac{M}{p} \Bigr)^2 .
  \label{eq:12.135}
\end{equation}
If we recall that we originally defined $\rho_m = m^2/M^2$, this is just a
complicated way of saying that, when $p$ becomes of order $m$, the mass term
becomes an important term in the Lagrangian.  At this point, the correlations
in the $\phi$ field begin to die away exponentially.  The characteristic range
of correlations, which in statistical mechanics would be called the
\emph{correlation length} $\xi$, is given by
\begin{equation}
  \xi \sim p_0^{-1},
  \qquad\text{where } \rho_m(p_0) = 1 .
  \label{eq:12.136}
\end{equation}
If we evaluate this criterion, we find $\xi \sim (M^2\rho_m)^{-1/2}$, that is,
$\xi \sim m^{-1}$, as we would have expected.

However, the application of this criterion at the fixed point $\lambda^{*}$
gives a much more interesting result.  If we set $\bar{\lambda} =
\lambda^{*}$, then Eq.~\eqref{eq:12.134} has the solution
\begin{equation}
  \rho_m \sim \rho_m(M)\ \Bigl( \frac{p}{M} \Bigr)^{-2 +
    \gamma_{\phi^2}(\lambda^{*})} .
  \label{eq:12.137}
\end{equation}
This gives a nontrivial relation
\begin{equation}
  \xi \sim \rho_m^{-\nu} ,
  \label{eq:12.138}
\end{equation}
where the exponent $\nu$ is given formally by the expression
\begin{equation}
  \nu^{-1} = 2 - \gamma_{\phi^2}(\lambda^{*}) .
  \label{eq:12.139}
\end{equation}
Using the results Eq.~(12.133) and \eqref{eq:12.117}, we can evaluate
this explicitly for $d$ near 4:
\begin{equation}
  \nu^{-1} = 2 - \frac{1}{6}(4-d) .
  \label{eq:12.140}
\end{equation}
Wilson and Fisher showed that this expression can be extended to a systematic
expansion of $\nu$ in powers of $\epsilon = (4-d)$.%
\footnote{This expansion is displayed to rather high order in E.~Brezin,
J.~C.~Le~Guillou, and J.~Zinn-Justin, \emph{Phys.\ Rev.} \textbf{D9}, 1121
(1974).}%
\footnote{K.~G.~Wilson and M.~E.~Fisher, \emph{Phys.\ Rev.\ Lett.} \textbf{28},
240 (1972).}

Because the exponent $\nu$ has an interpretation in statistical mechanics, it
is directly measurable in the realistic case of three dimensions.  In the
statistical mechanical interpretation of scalar field theory, $\rho_m$ is
just the parameter that one must adjust finely to bring the system to the
critical temperature.  Thus $\rho_m$ is proportional to the deviation from
the critical temperature, $(T-T_c)$.  Our field theoretic analysis thus
implies that the correlation length in a magnet grows as $T\to T_c$ according
to the scaling relation
\begin{equation}
  \xi \sim (T-T_c)^{-\nu} .
  \label{eq:12.141}
\end{equation}
It also gives a definite, and somewhat unusual, prediction for the value of
$\nu$.  It predicts that $\nu$ is close to the value $1/2$ suggested by the
Landau approximation studied in Chapter~8 (Eq.~\eqref{eq:8.16}), but that
$\nu$ differs from this value by some systematic corrections.

A scaling behavior of the type \eqref{eq:12.141} is observed in magnets, and
it is known that several definite scaling laws occur, depending on the
symmetry of the spin ordering.  Magnets can be characterized by the number of
fluctuating spin components: $N=1$ for magnets with a preferred axis, $N=2$
for magnets with a preferred plane, and $N=3$ for magnets that are isotropic
in three-dimensional space.  The experimental value of $\nu$ depends on this
parameter.  The $\phi^4$ field theory discussed in this chapter contained
only one fluctuating field; this is the analogue of a magnet with one spin
component.  In Chapter~11, we considered a generalization of $\phi^4$ theory
to a theory of $N$ fields with $O(N)$ symmetry.  We might guess that this
system models magnets of general $N$.

If this correspondence is correct, Eq.~\eqref{eq:12.140} gives a prediction
for the value of $\nu$ in magnets with a preferred axis.  In Section~13.1, we
will repeat the analysis leading to this equation in the $O(N)$-symmetric
$\phi^4$ theory and derive the formula
\begin{equation}
  \nu^{-1} = 2 - \frac{N+2}{N+8}(4-d) ,
  \label{eq:12.142}
\end{equation}
valid for general $N$ to first order in $(4-d)$.  For the cases $N = 1, 2, 3$
and $d = 3$, this formula predicts
\begin{equation}
  \nu = 0.60,\ 0.63,\ 0.65 .
  \label{eq:12.143}
\end{equation}
For comparison, the best current experimental determinations of $\nu$ in
magnetic systems give%
\footnote{For further details, see Table~13.1 and the accompanying
discussion.}
\begin{equation}
  \nu = 0.64,\ 0.67,\ 0.71
  \label{eq:12.144}
\end{equation}
for $N = 1, 2, 3$.  The prediction \eqref{eq:12.143} gives a reasonable first
approximation to the experimental results.

The ability of quantum field theory to predict the critical exponents gives a
concrete application both of the formal connection between quantum field
theory and statistical mechanics and of the flows of coupling constants
predicted by the renormalization group.  However, there is another
experimental aspect of critical behavior that is even more remarkable, and
more persuasive.  Critical behavior can be studied not only in magnets but
also in fluids, binary alloys, superfluid helium, and a host of other systems.
It has long been known that, for systems with this disparity of microscopic
dynamics, the scaling exponents at the critical point depend only on the
dimension $N$ of the fluctuating variables and not on any other detail of
the atomic structure.  Fluids, binary alloys, and uniaxial magnets, for
example, have the same critical exponents.  To the untutored eye, this seems
to be a miracle.  But for a quantum field theorist, this conclusion is the
natural outcome of the renormalization group idea, in which most details of
the field theoretic interaction are described by operators that become
irrelevant as the field theory finds its proper, simple, large-distance
behavior.

\problems

\textbf{12.1} \emph{Beta functions in Yukawa theory.}  In the pseudoscalar
Yukawa theory studied in Problem~10.2, with masses set to zero,
\begin{equation}
  \mathcal{L} = \tfrac{1}{2}(\partial_\mu\phi)^2 - \frac{\lambda}{4!}\phi^4
    - ig\bar{\psi}\gamma^5\psi\phi ,
\end{equation}
compute the Callan-Symanzik $\beta$ functions for $\lambda$ and $g$:
\begin{equation}
  \beta_\lambda(\lambda, g),\qquad \beta_g(\lambda, g) .
\end{equation}
[Hint: Compute the divergences of the scalar and fermion propagators and the
two types of vertex, and extract the coefficients of the divergent logarithms
as in Eqs.~\eqref{eq:12.50} and \eqref{eq:12.54}.]

\textbf{12.2} \emph{Critical exponents in the $O(N)$ model.}  Repeat the
analysis of Section~12.5 for the $O(N)$-symmetric $\phi^4$ theory, and derive
the formula \eqref{eq:12.142} for the critical exponent $\nu$ to first order
in $\epsilon = 4-d$.

\textbf{12.3} \emph{Evolution of the fine-structure constant.}  Using the QED
$\beta$ function, compute the value of the fine-structure constant at the
$Z$ pole, given its value at zero momentum.  Show that
$\alpha(m_Z) \approx 1/128$, as compared to $\alpha(0) \approx 1/137$.

\textbf{12.4} \emph{The anomalous dimension of the $\phi^2$ operator at two
loops.}  Compute the two-loop contribution to the anomalous dimension of the
mass operator in $\phi^4$ theory.
